\begin{derivation}[从\eqnrefs{eq:qcd-soft-leg-factor-dp68}到\eqnrefs{eq:qcd-color-charge-conservation-amplitude-dp10}]
先取一条出射无质量夸克外腿。把不含额外软胶子的其余硬过程记为颜色--旋量对象 $\mathcal H_i$，则把四动量 $k$、颜色 $a$、偏振 $\varepsilon_\lambda$ 的胶子接到该外腿后，相关因子为
\begin{equation*}
 \mathcal M_{i}^{a}(k,\lambda)
 =g_s\,\bar u(p_i)\slashed\varepsilon_\lambda^*(k)T_i^a
 \frac{\slashed p_i+\slashed k}{(p_i+k)^2+\ii0}\,\mathcal H_i .
\end{equation*}
外腿与软胶子都在壳，$p_i^2=k^2=0$，所以传播子分母为
\begin{equation*}
 (p_i+k)^2=2p_i\cdot k.
\end{equation*}
分子中先处理不含 $k$ 的部分。由 Clifford 代数
\begin{equation*}
 \slashed\varepsilon^*\slashed p_i
 =2p_i\cdot\varepsilon^*-\slashed p_i\slashed\varepsilon^*
\end{equation*}
以及外腿 Dirac 方程 $\bar u(p_i)\slashed p_i=0$，有
\begin{align*}
 \bar u(p_i)\slashed\varepsilon^*(\slashed p_i+\slashed k)
 &=2p_i\cdot\varepsilon^*\,\bar u(p_i)
   +\bar u(p_i)\slashed\varepsilon^*\slashed k,\\
 \frac{\bar u\slashed\varepsilon^*\slashed k}{2p_i\cdot k}
 &=\order(k^0),
 \qquad
 \frac{2p_i\cdot\varepsilon^*\,\bar u}{2p_i\cdot k}
 =\frac{p_i\cdot\varepsilon^*}{p_i\cdot k}\bar u.
\end{align*}
因此随软能量 $\omega$ 最强的 $\omega^{-1}$ 项为
\begin{equation*}
 \mathcal M_i^{a}(k,\lambda)
 =g_sT_i^a\frac{p_i\cdot\varepsilon_\lambda^*}{p_i\cdot k}\,
 \mathcal M_n+\order(\omega^0).
\end{equation*}
入射外腿通过 crossing 写成全出射 convention 时会多出相反的颜色荷方向。统一记
\begin{equation*}
 \sigma_i^{\rm io}=
 \begin{cases}
 +1,&\text{出射外腿},\\
 -1,&\text{入射外腿},
 \end{cases}
\end{equation*}
即可得到正文\eqnrefs{eq:qcd-soft-leg-factor-dp68}。软胶子在领先阶不能分辨硬子图内部的短距离结构，故把所有外腿贡献相加：
\begin{equation*}
 |\mathcal M_{n+1}^{a}(k,\lambda)\rangle
 =g_s\sum_i\sigma_i^{\rm io}\mathbf T_i^a
 \frac{p_i\cdot\varepsilon_\lambda^*}{p_i\cdot k}
 |\mathcal M_n\rangle+\order(\omega^0),
\end{equation*}
即正文\eqnrefs{eq:qcd-soft-gluon-theorem-dp10}。

颜色荷守恒不是附加假设，而由同一振幅的 Ward 恒等式给出。把外胶子偏振替换为纵向方向 $\varepsilon_\lambda^\mu\to k^\mu$，领先软项变成
\begin{align*}
 \left.|\mathcal M_{n+1}^{a}\rangle\right|_{\varepsilon\to k}
 &=g_s\sum_i\sigma_i^{\rm io}\mathbf T_i^a
 \frac{p_i\cdot k}{p_i\cdot k}|\mathcal M_n\rangle+\order(k)\\
 &=g_s\left(\sum_i\sigma_i^{\rm io}\mathbf T_i^a\right)|\mathcal M_n\rangle+\order(k).
\end{align*}
物理外胶子振幅满足 $k_\mu\mathcal M^{a\mu}=0$。取 $k\to0$ 后，$\order(k)$ 项消失，于是
\begin{equation*}
 \left(\sum_i\sigma_i^{\rm io}\mathbf T_i^a\right)|\mathcal M_n\rangle=0,
\end{equation*}
得到正文\eqnrefs{eq:qcd-color-charge-conservation-amplitude-dp10}。因此软定理中的纵向规范依赖与硬振幅的总颜色荷恰好由同一 Ward 恒等式消去。
\end{derivation}

\begin{derivation}[从\eqnrefs{eq:qcd-collinear-virtuality-dp68}到\eqnrefs{eq:q-to-qg-splitting-probability-dp10}]
令类光母方向 $P^\mu$ 与辅助类光方向 $n^\mu$ 满足
\begin{equation*}
 P^2=n^2=0,
 \qquad P\cdot n\neq0,
\end{equation*}
并取横向矢量 $k_\perp\cdot P=k_\perp\cdot n=0$。记其欧几里得平方为
\begin{equation*}
 K^2\equiv-k_\perp^2>0.
\end{equation*}
使两个子粒子严格在壳的 Sudakov 参数化为
\begin{align*}
 p^\mu
 &=zP^\mu+k_\perp^\mu
 +\frac{K^2}{z}\frac{n^\mu}{2P\cdot n},\\
 k^\mu
 &=(1-z)P^\mu-k_\perp^\mu
 +\frac{K^2}{1-z}\frac{n^\mu}{2P\cdot n}.
\end{align*}
直接平方可得
\begin{align*}
 p^2
 &=-K^2+2z(P\cdot n)\frac{K^2}{2z(P\cdot n)}=0,\\
 k^2
 &=-K^2+2(1-z)(P\cdot n)
 \frac{K^2}{2(1-z)(P\cdot n)}=0.
\end{align*}
离壳母动量 $\widetilde P=p+k$ 因而为
\begin{equation*}
 \widetilde P^\mu
 =P^\mu+\frac{K^2}{z(1-z)}\frac{n^\mu}{2P\cdot n},
\end{equation*}
所以
\begin{equation*}
 \widetilde P^2=\frac{K^2}{z(1-z)},
\end{equation*}
即正文\eqnrefs{eq:qcd-collinear-virtuality-dp68}。

把分裂顶角连到任意硬子图 $\mathcal H$。共线极限下，与硬子图相连的旋量密度矩阵只需保留母方向 $\slashed P$。对子夸克自旋和胶子的两个物理偏振求和时，先在轴向规范 $n\cdot A=0$ 中选择两个横向基 $e_a^\mu$，满足
\begin{equation*}
 e_a\cdot P=e_a\cdot n=0,
 \qquad e_a\cdot e_b=-\delta_{ab}.
\end{equation*}
写 $k_\perp^\mu=\kappa_a e_a^\mu$。满足 $k\cdot\varepsilon_a=n\cdot\varepsilon_a=0$ 的物理偏振可取
\begin{equation*}
 \varepsilon_a^\mu(k)
 =e_a^\mu-\frac{\kappa_a}{(1-z)P\cdot n}n^\mu.
\end{equation*}
分裂部分的自旋密度矩阵为
\begin{equation*}
 \mathcal S
 =\sum_{a=1}^{2}
 \frac{\slashed{\widetilde P}\slashed\varepsilon_a^*\slashed p
 \slashed\varepsilon_a\slashed{\widetilde P}}
 {(\widetilde P^2)^2}.
\end{equation*}
在领先共线幂次下写 $\mathcal S=A(z,K^2)\slashed P+\text{次领项}$。用 $\slashed n$ 投影即可确定 $A$：
\begin{equation*}
 4(P\cdot n)A
 =\operatorname{tr}(\slashed n\mathcal S).
\end{equation*}
这里的六伽马迹不能作为黑箱略过。两个物理偏振的完备关系为
\begin{equation*}
 \sum_{a=1}^{2}\varepsilon_a^{*\mu}\varepsilon_a^{\nu}
 =-\eta^{\mu\nu}
 +\frac{k^\mu n^\nu+n^\mu k^\nu}{k\cdot n},
 \qquad k\cdot n=(1-z)P\cdot n,
\end{equation*}
其中 $n^2=0$。因此先在两个偏振指标上完成收缩：
\begin{align*}
\sum_a\slashed\varepsilon_a^*\slashed p\slashed\varepsilon_a
={}&-\gamma_\mu\slashed p\gamma^\mu
 +\frac{\slashed k\slashed p\slashed n
 +\slashed n\slashed p\slashed k}{k\cdot n}\\
={}&2\slashed p
 +\frac{\slashed k\slashed p\slashed n
 +\slashed n\slashed p\slashed k}{(1-z)P\cdot n},
\end{align*}
其中用了四维 Clifford 恒等式
$\gamma_\mu\slashed p\gamma^\mu=-2\slashed p$。

于是所需迹分成三个可以逐项复核的标量：
\begin{align*}
\mathcal T_0
&\equiv2\Tr(\slashed n\slashed{\widetilde P}\slashed p\slashed{\widetilde P}),\\
\mathcal T_1
&\equiv\frac{1}{(1-z)P\cdot n}
\Tr(\slashed n\slashed{\widetilde P}\slashed k\slashed p
\slashed n\slashed{\widetilde P}),\\
\mathcal T_2
&\equiv\frac{1}{(1-z)P\cdot n}
\Tr(\slashed n\slashed{\widetilde P}\slashed n\slashed p
\slashed k\slashed{\widetilde P}).
\end{align*}
由 Sudakov 参数化可直接读出
\begin{equation*}
 n\cdot\widetilde P=P\cdot n,
 \quad n\cdot p=zP\cdot n,
 \quad n\cdot k=(1-z)P\cdot n,
 \quad \widetilde P^2\equiv Q^2=\frac{K^2}{z(1-z)},
\end{equation*}
以及
\begin{equation*}
 \widetilde P\cdot p
 =\widetilde P\cdot k
 =p\cdot k=\frac{Q^2}{2}.
\end{equation*}
先用四伽马迹
$\Tr(\slashed a\slashed b\slashed c\slashed d)
=4(a\!\cdot b\,c\!\cdot d-a\!\cdot c\,b\!\cdot d+a\!\cdot d\,b\!\cdot c)$，得到
\begin{align*}
\mathcal T_0
&=2\times4\left[
 (n\cdot\widetilde P)(p\cdot\widetilde P)
 -(n\cdot p)\widetilde P^2
 +(n\cdot\widetilde P)(\widetilde P\cdot p)
 \right]\\
&=8(P\cdot n)Q^2(1-z).
\end{align*}
六伽马迹也逐项展开。由 Clifford 反对易关系与迹循环性，任意六个四矢量满足递推式
\begin{align*}
\Tr(\slashed a\slashed b\slashed c\slashed d\slashed e\slashed f)
={}&(a\cdot b)\Tr(\slashed c\slashed d\slashed e\slashed f)
-(a\cdot c)\Tr(\slashed b\slashed d\slashed e\slashed f)\\
&+(a\cdot d)\Tr(\slashed b\slashed c\slashed e\slashed f)
-(a\cdot e)\Tr(\slashed b\slashed c\slashed d\slashed f)\\
&+(a\cdot f)\Tr(\slashed b\slashed c\slashed d\slashed e).
\end{align*}
对第一项的分子
\(\Tr(\slashed n\slashed{\widetilde P}\slashed k\slashed p
\slashed n\slashed{\widetilde P})\)，因为 $n^2=0$，递推后只剩
\begin{align*}
={}&(n\cdot\widetilde P)\Tr(\slashed k\slashed p\slashed n\slashed{\widetilde P})
-(n\cdot k)\Tr(\slashed{\widetilde P}\slashed p\slashed n\slashed{\widetilde P})\\
&+(n\cdot p)\Tr(\slashed{\widetilde P}\slashed k\slashed n\slashed{\widetilde P})
+(n\cdot\widetilde P)\Tr(\slashed{\widetilde P}\slashed k\slashed p\slashed n).
\end{align*}
四个四伽马迹依次为
\begin{align*}
\Tr(\slashed k\slashed p\slashed n\slashed{\widetilde P})&=4z(P\cdot n)Q^2,\\
\Tr(\slashed{\widetilde P}\slashed p\slashed n\slashed{\widetilde P})&=4z(P\cdot n)Q^2,\\
\Tr(\slashed{\widetilde P}\slashed k\slashed n\slashed{\widetilde P})&=4(1-z)(P\cdot n)Q^2,\\
\Tr(\slashed{\widetilde P}\slashed k\slashed p\slashed n)&=4z(P\cdot n)Q^2.
\end{align*}
把 $n\cdot k=(1-z)P\cdot n$ 与 $n\cdot p=zP\cdot n$ 代回，第二、第三项恰好抵消，故
\begin{equation*}
\Tr(\slashed n\slashed{\widetilde P}\slashed k\slashed p
\slashed n\slashed{\widetilde P})
=8z(P\cdot n)^2Q^2.
\end{equation*}
完全同样地，对
\(\Tr(\slashed n\slashed{\widetilde P}\slashed n\slashed p
\slashed k\slashed{\widetilde P})\) 递推，有
\begin{align*}
={}&(n\cdot\widetilde P)\Tr(\slashed n\slashed p\slashed k\slashed{\widetilde P})
+(n\cdot p)\Tr(\slashed{\widetilde P}\slashed n\slashed k\slashed{\widetilde P})\\
&-(n\cdot k)\Tr(\slashed{\widetilde P}\slashed n\slashed p\slashed{\widetilde P})
+(n\cdot\widetilde P)\Tr(\slashed{\widetilde P}\slashed n\slashed p\slashed k),
\end{align*}
其中四个迹分别为
\begin{equation*}
4z(P\cdot n)Q^2,
\quad4(1-z)(P\cdot n)Q^2,
\quad4z(P\cdot n)Q^2,
\quad4z(P\cdot n)Q^2.
\end{equation*}
中间两项再次抵消，所以该六伽马迹同样等于
\(8z(P\cdot n)^2Q^2\)。除以共同的
$(1-z)P\cdot n$ 后，
\begin{align*}
\mathcal T_1
&=\frac{8(P\cdot n)Q^2z}{1-z},\\
\mathcal T_2
&=\frac{8(P\cdot n)Q^2z}{1-z}.
\end{align*}
所以总迹不是直接引用的“标准结果”，而是
\begin{align*}
\mathcal T_0+\mathcal T_1+\mathcal T_2
&=8(P\cdot n)Q^2
\left[(1-z)+\frac{2z}{1-z}\right]\\
&=\frac{8(P\cdot n)Q^2}{1-z}(1+z^2)\\
&=\frac{8(P\cdot n)K^2}{z(1-z)^2}(1+z^2).
\end{align*}
这一步清楚显示了 $1+z^2$ 的来源：它是两个横向物理偏振与费米子 Clifford 代数共同留下的自旋结构，而 $1/(1-z)$ 的额外增强随后还会与共线相空间结合。

再用
\begin{equation*}
 (\widetilde P^2)^2
 =\frac{K^4}{z^2(1-z)^2},
\end{equation*}
得到
\begin{equation*}
 A(z,K^2)
 =\frac{2z(1+z^2)}{K^2}.
\end{equation*}
颜色平均给 $T^aT^a=C_F\mathbf1$，于是分裂部分对任意硬过程的平方振幅只通过普适因子
\begin{equation*}
 g_s^2C_F\frac{2z(1+z^2)}{K^2}
\end{equation*}
进入。

还需把 $(n+1)$ 体相空间在共线区因子化，而不能把相空间因子当作额外规则引用。对子系统总动量
\begin{equation*}
 Q^\mu\equiv p^\mu+k^\mu=\widetilde P^\mu,
 \qquad Q^2\equiv s_{pk}>0
\end{equation*}
插入恒等分解，Lorentz 不变相空间的递推式为
\begin{equation*}
 \dd\Phi_{n+1}
 =\dd\Phi_n(\ldots,Q)\,
 \frac{\dd s_{pk}}{2\pi}\,
 \dd\Phi_2(Q;p,k).
\end{equation*}
现在只需把二体子空间的角变量换成纵向动量分数 $z$。在 $Q$ 的瞬时静止系中，两条无质量子动量大小相同，
\begin{equation*}
 |\mathbf p_*|=|\mathbf k_*|=\frac{\sqrt{s_{pk}}}{2},
\end{equation*}
所以前文已经从 Lorentz 不变测度推得的二体相空间退化为
\begin{equation*}
 \dd\Phi_2(Q;p,k)
 =\frac{1}{32\pi^2}\,\dd\Omega_*
 =\frac{1}{32\pi^2}\,\dd\varphi_*\,\dd(\cos\theta_*).
\end{equation*}
选择辅助类光方向 $n^\mu$ 定义纵向轴。由于
\begin{equation*}
 z\equiv\frac{p\cdot n}{Q\cdot n},
\end{equation*}
在 $Q$ 静止系可取
\begin{equation*}
 n_*^\mu=\nu(1,0,0,-1),
 \qquad
 p_*^\mu=\frac{\sqrt{s_{pk}}}{2}
 (1,\sin\theta_*\cos\varphi_*,\sin\theta_*\sin\varphi_*,\cos\theta_*),
\end{equation*}
其中任意正比例因子 $\nu$ 会在比值中消去。于是
\begin{align*}
 p_*\cdot n_*
 &=\frac{\nu\sqrt{s_{pk}}}{2}(1+\cos\theta_*),\\
 Q_*\cdot n_*
 &=\nu\sqrt{s_{pk}},
\end{align*}
从而
\begin{equation*}
 z=\frac{1+\cos\theta_*}{2},
 \qquad
 \dd(\cos\theta_*)=2\,\dd z.
\end{equation*}
对未观测方位角积分 $\int_0^{2\pi}\dd\varphi_*$ 后，
\begin{align*}
 \int_0^{2\pi}\dd\varphi_*\,\dd\Phi_2
 &=\frac{1}{32\pi^2}(2\pi)(2\,\dd z)\\
 &=\frac{\dd z}{8\pi}.
\end{align*}
因此母虚度与二体子空间共同给
\begin{align*}
 \frac{\dd s_{pk}}{2\pi}\,\dd\Phi_2
 &\longrightarrow
 \frac{\dd s_{pk}\,\dd z}{16\pi^2}.
\end{align*}
最后使用已经由 Sudakov 参数化严格得到的
\begin{equation*}
 s_{pk}=\frac{K^2}{z(1-z)}.
\end{equation*}
在把 $(z,K^2)$ 选作局部相空间坐标后，固定 $z$ 取微分，
\begin{equation*}
 \left.\dd s_{pk}\right|_z
 =\frac{\dd K^2}{z(1-z)},
\end{equation*}
于是
\begin{equation*}
 \boxed{
 \dd\Phi_{\rm coll}
 =\frac{1}{16\pi^2}
 \frac{\dd z\,\dd K^2}{z(1-z)}.}
\end{equation*}
这一步也说明了 $1/[z(1-z)]$ 的来源：它不是分裂顶角的动力学奇点，而是把二体角变量与母虚度改写成 $(z,K^2)$ 时产生的 Jacobian。

因此分裂概率相对于硬过程为
\begin{align*}
 \dd P_{q\to qg}
 &=g_s^2C_F\frac{2z(1+z^2)}{K^2}
 \frac{1}{16\pi^2}
 \frac{\dd z\,\dd K^2}{z(1-z)}\\
 &=\frac{\alpha_s}{2\pi}
 C_F\frac{1+z^2}{1-z}
 \frac{\dd K^2}{K^2}\,\dd z,
\end{align*}
其中用了 $g_s^2=4\pi\alpha_s$。把 $K^2$ 重新记作正文使用的 $k_\perp^2>0$，即得\eqnrefs{eq:q-to-qg-splitting-probability-dp10}。
\end{derivation}
